\begin{note}
    Если $(X, \|\cdot\|)$ -- бесконечномерное нормированное пространство, и
    \[
        S^* = \left\{f \in X^* \mid \|f\| = 1\right\},
    \]
    то *-слабое замыкание $S^*$ является единичным нормированным шаром $B_1^*[0]$ в $X^*$.

    Пусть $g \in X^* \setminus B_1^*[0]$ (то есть $\|g\| > 1$).
    По определению нормы функционала
    \[
        \sup_{\|x\| = 1} |g(x)| > 1,
    \]
    а значит $\exists x_0 \in X\colon \|x_0\| = 1$ и $|g(x_0)| > 1$.
    Тогда $V(g, x_0, |g(x_0)| - 1) \cap B_1^*[0] = \emptyset$, так как
    \[
        \forall f \in V(g, x_0, |g(x_0)| - 1) \hookrightarrow |g(x_0)| - |f(x_0)| \leq |f(x_0) - g(x_0)| < |g(x_0)| - 1 \Rightarrow 1 < |f(x_0)| \leq \|f\|.
    \]
    И тогда дополнение $B_1^*[0]$ является *-слабо открытым.

    Аналогично случаю со слабой топологией, $B_1(0)$ вложено в замыкание, так как всякая *-слабая окрестность точки содержит бесконечномерное сдвинутое на точку подпространство (а значит отрезок, один конец которого имеет норму меньше единицы, а второй -- сколь угодно большую).
    Поэтому $[S^*]_{\tau_{w*}} = B_1^*[0]$.
\end{note}

\section{Неметризуемость слабых топологий.}\label{seq:unmetrizability}
\begin{theorem}
    Пусть $X$ -- бесконечномерное нормированное пространство.
    \begin{enumerate}
        \item Слабая топология в $X$ неметризуема.
        \item Если $X$ -- банахово, то *-слабая топология в $X^*$ неметризуема.
    \end{enumerate}
\end{theorem}
\begin{proof}
    Покажем неметризуемость слабой топологии.
    Пусть $\exists \rho$ -- метрика на $X$ такая, что $\tau_{\rho} = \tau_w$.
    Рассмотрим $O_r^\rho(0) = \{x \in X \mid \rho(x, 0) < r\}$.
    Известно, что $\{O_{1/n}^{\rho}(0)\}_{n = 1}^{\infty}$ образуют счётную локальную базу $\tau_\rho$.
    В частности, так как $O_{1/n}^{\rho}(0)$ слабо открыто, то $\exists f_{1, n}, \ldots, f_{m_n, n} \in X^*$ и $\epsilon_n > 0$ такие, что
    \[
        \bigcap_{k = 1}^{m_n} \ker f_{k, n} \subset \bigcap_{k = 1}^{m_n} V(0, f_{k, n}, \epsilon_n) \subset O^{\rho}_{1/n}(0).
    \]
    Но $\forall U(0) \in \tau_w = \tau_\rho$ существует $n$ такое, что
    \[
        O_{1/n}^{\rho}(0) \subset U(0).
    \]
    Теперь возьмём для $f \in X^*$ окрестность нуля как $V(0, f, 1) \supset O_{1/n}^{\rho}(0) \supset \bigcap_{k = 1}^{m_n} \ker f_{k, n}$.
    Аналогично лемме о сопряжённом пространстве $\ker f \supset \bigcap_{k = 1}^{m_n} \ker f_{k, n}$, а значит, вследствие леммы о ядрах $f$ лежит в линейной оболочке $\{f_{k, n}\}_{k = 1}^{m_n}$.
    Пусть $L_n = \Lin\{f_{k, n}\}_{k = 1}^{m_n}$.
    Тогда
    \[
        X^* = \bigcup_{n = 1}^{\infty} L_n.
    \]
    Поскольку все $L_n$ -- сильно замкнутые подпространства (так как конечномерны), то в силу банаховости $X^*$ одно из них имеет внутреннюю точку, а значит совпадает со всем пространством.
    Но тогда размерность сопряжённого пространства конечна -- противоречие.

    Покажем неметризуемость *-слабой топологии на $X^*$, если $X$ -- банахово.
    Пусть $\exists \rho^*$ -- метрика на $X^*$ такая, что $\tau_{\rho^*} = \tau_{w*}$.
    Рассмотрим $O_r^{\rho^*}(0) = \{f \in X^* \mid \rho^*(f, 0) < r\}$.
    Известно, что $\{O_{1/n}^{\rho^*}(0)\}_{n = 1}^{\infty}$ образуют счётную локальную базу $\tau_{\rho^*}$.
    В частности, так как $O_{1/n}^{\rho^*}(0)$ *-слабо открыто, то $\exists x_{1, n}, \ldots, x_{m_n, n} \in X$ и $\epsilon_n > 0$ такие, что
    \[
        \bigcap_{k = 1}^{m_n} \ker \Phi_{x_{k, n}} \subset \bigcap_{k = 1}^{m_n} V(0, x_{k, n}, \epsilon_n) \subset O^{\rho^*}_{1/n}(0).
    \]
    Но $\forall U(0) \in \tau_{w*} = \tau_{\rho^*}$ существует $n$ такое, что
    \[
        O_{1/n}^{\rho^*}(0) \subset U(0).
    \]
    Теперь возьмём для $x \in X$ окрестность нуля как $V(0, x, 1) \supset O_{1/n}^{\rho^*}(0) \supset \bigcap_{k = 1}^{m_n} \ker \Phi_{x_{k, n}}$.
    Аналогично лемме о сопряжённом пространстве $\ker \Phi_x \supset \bigcap_{k = 1}^{m_n} \ker \Phi_{x_{k, n}}$, а значит, вследствие леммы о ядрах $x$ лежит в линейной оболочке $\{x_{k, n}\}_{k = 1}^{m_n}$.
    Пусть $L_n = \Lin\{x_{k, n}\}_{k = 1}^{m_n}$.
    Тогда
    \[
        X = \bigcup_{n = 1}^{\infty} L_n.
    \]
    Поскольку все $L_n$ -- сильно замкнутые подпространства (так как конечномерны), то в силу банаховости $X$ одно из них имеет внутреннюю точку, а значит совпадает со всем пространством.
    Но тогда размерность пространства $X$ конечна -- противоречие.
\end{proof}